% MOMENTBRIDGE1
\documentclass[11pt]{article}
\usepackage[margin=.88in]{geometry}
\usepackage{amsmath,amssymb,amsthm}
\newtheorem{proposition}{Proposition}
\newtheorem{lemma}{Lemma}
\begin{document}

\begin{center}
{\Large MOMENTBRIDGE1: Moment-balanced UNSAT versus the uniform slice}\\[4pt]
\textbf{Date: 2026-10-04}\\
\textbf{Status: PROVED/EXACT parameter obstruction + COMPUTATIONAL tests}\\
\textbf{P versus NP status: OPEN}
\end{center}

\section{Why the next step changed}

The immediate question is no longer merely whether pair-matched rare UNSAT
formulas have a visible cubic statistic.  The literature gives a much more
specific target.

Wegener's slice-function theory identifies monotone circuits on a slice with
set-circuit structure and explicitly warns that the standard monotone
lower-bound methods do not automatically transfer to slice functions.
Cavalar--Kumar--Rossman (CKR) nevertheless show that slice sunflowers can be
used inside Razborov's approximation method for carefully chosen test
distributions.  Their construction uses a closure operator, trimming, and a
slice-sunflower bound.

Very recently, Cavalar--Fabris--Mukhopadhyay--Srinivasan--Yehudayoff
(ECCC TR26-220, 2026) proved robust sunflower bounds by a method-of-moments
argument for locally pseudo-independent distributions.  Their theorem is
particularly relevant because it needs only
\[
k=O(w\log(1/\delta))
\]
local order for a width-\(w\) monotone DNF, rather than a full PRG for DNFs.
The same paper emphasizes that its argument can apply even when the
distribution does not fool arbitrary DNFs.

Amireddy--Behera--Srinivasan--Sudan--Willumsgaard (ECCC TR26-152, 2026)
give low-degree tests directly on Boolean slices.  This confirms that
low-degree structure on the Johnson slice is robustly testable, but it does
not imply that general monotone circuits have low degree.

\section{A direct parameter obstruction to the off-the-shelf CKR template}

Let the ambient number of clause coordinates be
\[
M=8\binom q3,
\]
and let the positive layer have weight \(t=\Theta(q)\).  Suppose one tries to
use the CKR slice-sunflower closure argument with a negative test distribution
that is uniform on an \(m\)-slice, where \(m\le t\).

For an \(\ell\)-minterm of an \(\varepsilon\)-closed approximator, CKR's
slice-sunflower bound has the form
\[
|M_\ell(h)|
\le
\left(
B x\log x
\right)^\ell,
\qquad
x=
\frac{\log(\ell/\varepsilon)}{p},
\qquad
p=\frac mM.
\tag{1}
\]

If the positive distribution were the uniform exact-\(t\) slice, a fixed
\(\ell\)-set is contained with probability
\[
\frac{(t)_\ell}{(M)_\ell}.
\tag{2}
\]

Thus the standard union-bound estimate for the positive error contains
\[
R_\ell
=
\frac{(t)_\ell}{(M)_\ell}
\left(
B x\log x
\right)^\ell.
\tag{3}
\]

\begin{proposition}[DENSITY-CANCEL1]
Fix \(\ell\ge1\) and any error sequence with
\(\log(\ell/\varepsilon)\ge c_0>0\).
If \(m\le t=\Theta(q)\) and \(M=\Theta(q^3)\), then
\[
R_\ell^{1/\ell}
=
\Omega\left(
B\frac tm
\log(\ell/\varepsilon)
\log\left(
\frac{M}{m}\log(\ell/\varepsilon)
\right)
\right),
\tag{4}
\]
up to the usual fixed-\(\ell\) falling-factorial correction.  In particular,
the right-hand side diverges like \(\Omega(\log q)\) for fixed
\(\ell\) and fixed positive \(B\).  Hence this direct CKR
closure-plus-union-bound parameterization cannot make the positive error
vanish on two nearby sparse slices.
\end{proposition}

\begin{proof}
Insert \(p=m/M\) into (1) and multiply by the \(\ell\)-th root of
(2).  For fixed \(\ell\),
\[
\left(\frac{(t)_\ell}{(M)_\ell}\right)^{1/\ell}
=
\frac tM(1+o(1)).
\]
Therefore
\[
R_\ell^{1/\ell}
=
B\frac tm
\log(\ell/\varepsilon)
\log\left(
\frac{M}{m}\log(\ell/\varepsilon)
\right)
(1+o(1)).
\]
Since \(t/m\ge1\) and \(M/m=\Theta(q^2)\), the logarithmic factor
diverges.
\end{proof}

This is a no-go for the \textbf{direct proof template}, not for all
sunflower or approximation arguments.

The 2026 pseudo-independent sunflower theorem does not automatically remove
this density cancellation: its generic threshold contains a factor of roughly
\((w/q_0)\log(1/\delta))^w\).  Multiplying by a same-density positive
containment probability leaves a factor of order
\((w\log(1/\delta))^w\).  A successful argument must therefore use more than
the old minterm-count plus union-bound step.

\section{A better distribution pair}

For circuit lower bounds, the clean negative distribution is the
\textbf{uniform exact-\(t\) slice itself}.  At density 3.5 this distribution
is SAT with high probability by the rigorous random-3SAT lower bound, so an
exact UNSAT circuit is 0 on almost all of it.

Instead of reweighting SAT formulas, define a symmetry-invariant positive
family over UNSAT formulas:
\[
\mu_{\lambda,q}(F)
\propto
\mathbf 1[
|F|=t\text{ and }F\text{ UNSAT}
]
\exp\left(
\lambda\cdot Z(F)
\right),
\tag{5}
\]
where \(Z(F)\) is the vector of the eight nontrivial signed pair-orbit
counts.

Because \(Z\) is invariant under signed variable permutations, every
\(\mu_{\lambda,q}\) is also invariant.  In particular, its one-coordinate
marginals remain exactly \(t/M\).

The target pair vector for the uniform exact slice is known exactly.  For a
pair orbit \((s,h)\) with \(n_{s,h}\) possible partners of a fixed clause,
\[
\mathbb E_{unif} P_{s,h}
=
n_{s,h}
\frac{t(t-1)}{2(M-1)}.
\tag{6}
\]

\section{POS-TILT1: empirical entropy cost of exact pair balancing}

The maximum-entropy exponential tilt was fit to empirical UNSAT samples so
that the weighted pair vector equals (6) to numerical precision.

Using all available UNSAT samples in each pooled chunk gave:

\[
\begin{array}{c|c|c|c|c|c}
q&\#U&ESS&ESS/\#U&D_{KL}(w\|\text{uniform sample})&
\|\lambda\|_2\\
\hline
16&324&35.55&0.110&1.140&1.644\\
18&322&32.66&0.101&1.453&2.241\\
20&976&116.81&0.120&1.246&1.828\\
24&827&83.47&0.101&1.341&1.886
\end{array}
\]

All four maximum-entropy fits reached pair-moment error below about
\(10^{-10}\).

\textbf{COMPUTATIONAL interpretation:} over \(q=16\)--\(24\), exact empirical
pair balancing costs about 1--1.5 nats of relative entropy and leaves roughly
10--12 percent effective sample mass.  There is no visible growth with
\(q\) in this small range.

This does not prove bounded asymptotic KL divergence or even asymptotic
existence of the exact tilt.

Train/holdout fitting is substantially noisier because rare UNSAT sampling is
itself a tail event.  For example, a single \(q=24\) chunk can fail empirical
convex-hull feasibility, while pooling three independent chunks restores it.

\section{Third-order hostile validation}

A canonical three-clause orbit that appeared negative in several earlier
pair-balanced runs was frozen before two new exact-pair-fit pools were tested.
Its incidence signature is
\[
((-,-,-),(-,-,0),(-,0,0),(0,-,+),(0,0,-)).
\tag{7}
\]
Equivalently, one variable is common to all three clauses with the same sign,
one additional variable is shared by the first two clauses with the same
sign, and another variable is shared by the second and third clauses with
opposite signs.

Two independent \(q=20\), \(N=14000\) pools were pair-balanced exactly.
They had 311 and 321 UNSAT samples, with ESS fractions 0.178 and 0.205.

The predeclared orbit effect, measured in pooled-formula standard deviations,
was
\[
-0.251
\qquad\text{and}\qquad
-0.017.
\]

Thus the apparent orbit does \textbf{not} robustly replicate after exact
pair balancing.

\section{Random spread-DNF probe}

The 2026 method-of-moments sunflower theorem motivates testing spread DNFs
directly rather than only individual orbit counts.

For each independent exact-pair-balanced \(q=20\) UNSAT pool, random monotone
DNFs over clause coordinates were generated independently of the formulas.

A width-2 DNF with 10,000 random terms gave acceptance gaps
\[
-0.0155
\qquad\text{and}\qquad
+0.0043
\]
between pair-balanced UNSAT and the uniform exact slice.

For width 3, ten independent random DNFs with 300,000 terms gave gaps
\[
\begin{split}
&0.1308,-0.0390,-0.0162,0.0462,-0.0547,\\
&-0.0174,0.0196,0.0168,-0.0432,0.0602.
\end{split}
\tag{8}
\]
Their mean gap was
\[
0.0103,
\]
with sample standard deviation
\[
0.0572.
\]

The isolated 0.1308 gap does not replicate.

\textbf{COMPUTATIONAL interpretation:} random spread width-2 and width-3 DNFs
do not show a stable systematic advantage against the pair-balanced UNSAT
distribution in these finite tests.

This is not a theorem for all spread DNFs, let alone all monotone circuits.

\section{What the literature now suggests}

The old target
\[
\text{small monotone circuit}
\Longrightarrow
\text{low harmonic degree}
\]
is probably too strong and is not supplied by slice harmonic analysis.

A more literature-aligned target is:

\[
\boxed{
\text{small monotone circuit}
\Longrightarrow
\text{two-distribution slice approximator built from spread bounded-width DNFs}
}
\tag{9}
\]

where the approximation is measured only on

\[
\mu_+ = \text{pair-/moment-balanced UNSAT}
\]
and
\[
\mu_- = \text{uniform exact }t\text{-slice}.
\]

The 2026 method-of-moments sunflower argument is relevant because it controls
spread DNFs using only local moments and dependency patterns.  However, its
published theorem is not yet the slice-intrinsic two-distribution circuit
approximation theorem required in (9).

Wegener's set-circuit formulation is therefore again important: any successful
argument should remain on the exact slice rather than relying on the standard
off-slice closure operation, which becomes degenerate for a Valiant slice
extension.

\section{Highest-value next theorem}

\textbf{TWO-DIST-MOMENT-APPROX.}  Construct or falsify an exact-slice
approximation scheme with the following shape.

For a monotone circuit \(C\) of size \(s\), build an approximator \(A_C\)
such that
\[
\Pr_{\mu_+}[C\ne A_C]
+
\Pr_{\mu_-}[C\ne A_C]
\le s\eta,
\tag{10}
\]
and every approximator is assembled from spread bounded-width DNFs whose
acceptance probabilities are determined, up to \(o(1)\), by moments through
order
\[
k=O(w\log(1/\eta)).
\tag{11}
\]

Then attempt to choose a positive UNSAT distribution whose moments through
that order match the uniform slice closely enough that every approximator
makes a constant error.

Mandatory hostile control: XOR-SLICE-EASY.  Any approximation statement based
only on generic slice moment matching that would also apply unchanged to
Gaussian-eliminable 3XOR slice functions is too strong and must be rejected.

\section{Status}

DENSITY-CANCEL1: \textbf{PROVED/EXACT} as a parameter obstruction to the
direct CKR union-bound template.

POS-TILT1, triple-orbit validation, and spread-DNF probes:
\textbf{COMPUTATIONAL}.

Asymptotic bounded-KL pair balancing: \textbf{OPEN}.

TWO-DIST-MOMENT-APPROX: \textbf{OPEN}.

General monotone slice lower bound and \(P\ne NP\): \textbf{OPEN}.

\begin{thebibliography}{7}
\bibitem{Wegener1985}
I. Wegener, On the Complexity of Slice Functions,
Theoretical Computer Science 38:55--68, 1985.

\bibitem{CKR2022}
B. P. Cavalar, M. Kumar, and B. Rossman,
Monotone Circuit Lower Bounds from Robust Sunflowers,
Algorithmica 84:3655--3685, 2022.

\bibitem{Rao2020}
A. Rao, Coding for Sunflowers,
Discrete Analysis, 2020.

\bibitem{CFMSY2026}
B. P. Cavalar, T. Fabris, P. Mukhopadhyay, S. Srinivasan, and A. Yehudayoff,
Derandomized Sunflowers and Radical Lower Bounds,
ECCC TR26-220, 2026.

\bibitem{ABSSW2026}
P. Amireddy, A. R. Behera, S. Srinivasan, M. Sudan, and S. V. Willumsgaard,
Low-Degree Testing Over Boolean Slices,
ECCC TR26-152, 2026.

\bibitem{FilmusMossel2016}
Y. Filmus and E. Mossel,
Harmonicity and Invariance on Slices of the Boolean Cube,
CCC 2016.
\end{thebibliography}

\end{document}
